\documentclass[conference,10pt]{IEEEtran}
\usepackage[utf8]{inputenc}
\usepackage{amsmath,amssymb,amsfonts}
\usepackage{algorithmic}
\usepackage{graphicx}
\usepackage{textcomp}
\usepackage{xcolor}
\usepackage{booktabs}
\usepackage{multirow}
\usepackage{url}
\usepackage{cite}
\usepackage{microtype}

\graphicspath{{figures/}}

\begin{document}

\title{Leakage-Aware, Recall-Constrained Evaluation of Tabular Models for Cardiovascular Risk Screening: A Multi-Seed Reimplementation and Extension of EnsCVDD-Net}

\author{
\IEEEauthorblockN{Person 1\IEEEauthorrefmark{1}, Person 2\IEEEauthorrefmark{2}, Person 3\IEEEauthorrefmark{3}, Person 4\IEEEauthorrefmark{4}}
\IEEEauthorblockA{Department of Computer Science and Engineering\\
Cardiovascular AI Research Laboratory\\
Email: \{\IEEEauthorrefmark{1}person1, \IEEEauthorrefmark{2}person2, \IEEEauthorrefmark{3}person3, \IEEEauthorrefmark{4}person4\}@university.edu}
}

\maketitle

\begin{abstract}
Recent literature has proposed hybrid deep neural architectures combining 1D Convolutional Neural Networks and Gated Recurrent Units (EnsCVDD-Net, BlCVDD-Net) for cardiovascular disease (CVD) detection from tabular survey data, reporting accuracies exceeding $91\%$ and precisions near $96\%$ under ADASYN oversampling. In this study, we conduct an auditable reimplementation and methodological audit of these architectures on the 253,680-record CDC BRFSS 2015 cohort. We mathematically prove that the reported precision/recall metrics are mutually incompatible with published accuracy under the natural positive survey prevalence ($9.42\%$), demonstrating that published benchmarks reflect an artificially balanced ($\approx 52.3\%$) evaluation distribution. We introduce a leakage-safe protocol that strictly confines preprocessing, feature ranking, and class weighting to training folds, preserving untouched natural prevalence on validation and test partitions. To address the invalid inductive bias of sequential 1D convolutions on arbitrary column orderings, we formulate the \textbf{Feature-Group Gated Network (FGGN)} and an un-gated \textbf{UniformGroupNet}, structuring predictors into four physiologically grounded domain representations (\textit{Clinical History}, \textit{Lifestyle}, \textit{Self-Reported Health}, \textit{Demographics/Access}). Across a rigorous 5-seed study (951,300 persisted test predictions), Histogram Gradient Boosting achieves the top overall ranking performance (PR-AUC: $0.3717 \pm 0.0067$, ROC-AUC: $0.8483 \pm 0.0026$). While domain-grouped architectures significantly outperform EnsCVDD-Net (PR-AUC: $0.3563 \pm 0.0055$, paired $\Delta = +0.0126$, $p = 0.0011$), learned softmax gating provides no ranking benefit over uniform semantic group pooling ($0.3689$ vs. $0.3693$, $p = 0.448$), establishing an honest neutral finding regarding model complexity. An 80/20 probability ensemble of HistGBT and FGGN achieves superior ranking ($0.3728 \pm 0.0064$). Validation-fitted Isotonic regression restores probability calibration, reducing test Brier score from $0.1654$ to $0.0703$ (beating the $0.0853$ constant-prevalence baseline). Finally, demographic subgroup analysis reveals that precision scales directly with baseline stratum prevalence ($17.7\%$ in young adults vs. $25.7\%$ in older adults), underscoring that population prevalence dynamics must govern the operational deployment of automated screening models.
\end{abstract}

\begin{IEEEkeywords}
Cardiovascular disease screening, tabular deep learning, EnsCVDD-Net, class imbalance, probability calibration, decision curve analysis, data leakage.
\end{IEEEkeywords}

\section{Introduction}
Cardiovascular disease (CVD) remains the leading cause of global mortality, accounting for approximately 17.9 million deaths annually. In primary healthcare systems, population-scale risk assessment relies heavily on standardized, low-cost behavioral surveillance questionnaires such as the Behavioral Risk Factor Surveillance System (BRFSS) conducted by the Centers for Disease Control and Prevention (CDC). Unlike clinical cohorts containing laboratory assays or continuous biomarkers, national health surveys collect self-reported, low-cardinality discrete responses regarding pre-existing diagnoses, demographic attributes, and lifestyle factors.

Designing automated machine learning screening tools from survey data presents two fundamental challenges:
\begin{enumerate}
    \item \textbf{Extreme Class Imbalance:} In the standardized CDC BRFSS 2015 cohort, positive CVD prevalence is approximately $9.42\%$. In this regime, unconstrained accuracy is severely misleading; a trivial majority-class baseline that predicts all individuals are healthy achieves $90.58\%$ accuracy while providing zero screening sensitivity ($0\%$ recall).
    \item \textbf{Clinical Asymmetry \& Recall Priority:} Missing an individual with cardiovascular pathology (False Negative) carries devastating medical consequences, whereas a false alarm (False Positive) prompts non-invasive confirmatory testing. Consequently, clinical screening models must be deployed under explicit sensitivity constraints ($\text{Recall} \ge 80\%$) rather than default $0.50$ probability thresholds.
\end{enumerate}

Recently, Khan et al.~\cite{khan2024heart} proposed \textbf{EnsCVDD-Net} (a 1D CNN + GRU hybrid) and \textbf{BlCVDD-Net} (an MLP meta-blender) on the CDC BRFSS 2015 dataset, reporting $>91.7\%$ accuracy, $96.0\%$ precision, and $87.8\%$ recall under ADASYN oversampling. However, sequential convolutions on static survey columns impose questionable spatial and temporal inductive biases. Crucially, the mathematical consistency of reporting $96.0\%$ precision alongside $87.8\%$ recall and $91.7\%$ accuracy under $9.42\%$ natural prevalence warrants formal audit.

In this work, we present an auditable, multi-seed reimplementation, mathematical audit, and physiological domain-grouped extension of EnsCVDD-Net.

\section{Dataset Contract \& Data Integrity Audit}
\subsection{Cryptographic Fingerprint \& Contract}
All experiments strictly adhere to the Kaggle-derived CDC BRFSS 2015 Heart Disease Health Indicators benchmark:
\begin{itemize}
    \item \textbf{Dimensions:} 253,680 individual patient records across 21 predictor variables and 1 binary target (\texttt{HeartDiseaseorAttack}).
    \item \textbf{Cryptographic Hash:} Audited and cryptographically locked with SHA-256 hash \texttt{953fb23bb131736d7879539a071d3f6d1eaae8d4867ab05335d0ba89b94d5f99}.
    \item \textbf{Class Distribution:} 229,787 negatives ($90.58\%$) vs. 23,893 positives ($9.42\%$) with zero missing values.
\end{itemize}

\subsection{Deduplication Sensitivity Audit}
Our audit identified 23,899 duplicate records ($9.42\%$) within the dataset. In low-cardinality survey data, identical feature vectors represent independent respondents sharing identical risk profiles rather than clerical errors. Removing all duplicates shifts prevalence marginally to $10.32\%$, while model performance rankings remain identical (HistGradientBoosting achieves PR-AUC $0.3732$ raw vs. $0.3835$ deduplicated). To preserve comparability with Khan et al.~\cite{khan2024heart}, the full 253,680 cohort is maintained as the primary benchmark.

\subsection{Leakage-Safe Stratification Boundary}
We enforce a strict boundary: raw data is partitioned into stratified $70\%$ Training ($N=177,576$), $15\%$ Validation ($N=38,052$), and $15\%$ Test ($N=38,052$) sets. Positive prevalence is strictly preserved at $9.42\%$ across all partitions. Scalers, Point-Biserial Correlation (PBCC) rankings, and thresholds are fitted exclusively on training data.

\section{Related Work}
We contextualize our study across four pillars:
\begin{enumerate}
    \item \textbf{Tree Ensembles vs. Deep Learning:} Grinsztajn et al.~\cite{grinsztajn2022tree} and Shwartz-Ziv \& Armon~\cite{shwartz2022tabular} established that gradient-boosted decision trees consistently outperform neural networks on tabular benchmarks.
    \item \textbf{Feature Gating:} Dynamic feature selection via TabNet~\cite{arik2021tabnet} and Temporal Fusion Transformers~\cite{lim2021temporal} inspires our domain gating, but we operate over physiologically defined groups rather than arbitrary single columns.
    \item \textbf{Probability Calibration:} Van Calster et al.~\cite{vancalster2016calibration} and Wallace \& Dahabreh~\cite{wallace2012cost} formalized clinical calibration requirements, showing that cost-sensitive adjustments distort posterior probabilities.
    \item \textbf{Resampled Inference:} Nadeau \& Bengio~\cite{nadeau2003inference} formulated variance corrections for repeated holdouts sharing training overlap.
\end{enumerate}

\section{Methodological Audit of Published Work}
\subsection{The Prevalence Inconsistency Proof}
Analytically, Classification Accuracy $A$, Positive Prevalence $\pi$, Sensitivity (Recall $R$), and Positive Predictive Value (Precision $P$) are linked by:
\begin{equation}
A = 1 - \pi + \pi R \left(2 - \frac{1}{P}\right)
\label{eq:consistency}
\end{equation}

\textbf{Proof 1 (Natural Prevalence):} Substituting $\pi = 0.094186$, $R = 0.878$, and $P = 0.960$ into Eq.~(\ref{eq:consistency}) mathematically requires an Accuracy of:
\begin{equation}
A_{\text{derived}} = 98.51\%
\end{equation}
This contradicts the published accuracy of $91.70\%$ by $+6.81$ percentage points.

\textbf{Proof 2 (Precision Upper Bound):} Given $A = 0.917$, $R = 0.878$, and $\pi = 0.0942$, the theoretical maximum attainable precision is $53.63\%$, which is $42.37$ percentage points below the reported $96.0\%$.

\textbf{Proof 3 (Reconciled Prevalence):} Solving for the synthetic positive prevalence $\pi^*$ that reconciles all three reported metrics yields:
\begin{equation}
\pi^* = \frac{1 - 0.917}{1 - 0.878(2 - 1/0.96)} = 52.34\%
\end{equation}
This demonstrates that the reported numbers are mathematically characteristic of an artificially balanced ($\approx 50:50$) test partition.

\subsection{Empirical Failure of ADASYN on Survey Features}
ADASYN continuous synthetic interpolation creates fractional levels for discrete indicators (e.g., \texttt{Smoker}=0.63). On HistGradientBoosting, ADASYN drops PR-AUC from $0.3732$ to $0.3306$ ($-0.0426$), whereas Random Oversampling preserves the discrete domain ($0.3735$).

\section{Domain-Grouped Neural Architectures}
\subsection{Physiological Domain Grouping}
We partition the 21 predictors into four clinical domain groups:
\begin{enumerate}
    \item \textit{Clinical History} ($d_1=6$): \texttt{HighBP}, \texttt{HighChol}, \texttt{CholCheck}, \texttt{Stroke}, \texttt{Diabetes}, \texttt{DiffWalk}.
    \item \textit{Lifestyle Factors} ($d_2=6$): \texttt{BMI}, \texttt{Smoker}, \texttt{PhysActivity}, \texttt{Fruits}, \texttt{Veggies}, \texttt{HvyAlcoholConsump}.
    \item \textit{Self-Reported Health} ($d_3=3$): \texttt{GenHlth}, \texttt{MentHlth}, \texttt{PhysHlth}.
    \item \textit{Access \& Demographics} ($d_4=6$): \texttt{AnyHealthcare}, \texttt{NoDocbcCost}, \texttt{Sex}, \texttt{Age}, \texttt{Education}, \texttt{Income}.
\end{enumerate}

\subsection{Architectural Specifications}
\textbf{Feature-Group Gated Network (FGGN):}
Each domain $g$ is encoded by an MLP $E_g: \mathbb{R}^{d_g} \to \mathbb{R}^{16}$. Embeddings $\mathbf{z}_g = E_g(\mathbf{x}_g)$ are concatenated into $\mathbf{z}_{\text{joint}} \in \mathbb{R}^{64}$. A softmax gating head computes:
\begin{equation}
\mathbf{w} = \text{softmax}(W_{\text{gate}} \mathbf{z}_{\text{joint}} + b_{\text{gate}})
\end{equation}
The fused representation $\mathbf{z}_{\text{fused}} = \sum_{g=1}^4 w_g \mathbf{z}_g$ is passed to classifier $C(\mathbf{z}_{\text{fused}})$.

\textbf{UniformGroupNet (Ablation):} Replaces learned gating with equal $1/4$ uniform pooling: $\mathbf{z}_{\text{fused}} = \frac{1}{4} \sum_{g=1}^4 \mathbf{z}_g$.

\textbf{InstanceGatedNet (Extension):} Employs a non-linear bottleneck gate $\mathbf{w}_i = \text{softmax}(W_2 \text{ReLU}(W_1 \mathbf{z}_{\text{joint}, i} + b_1) + b_2)$ providing sample-specific dynamic routing.

\subsection{Operational Threshold Policy}
For every model, decision threshold $\tau^*$ is selected on the validation set to satisfy:
\begin{equation}
\tau^* = \max \{ \tau \in [0, 1] \mid \text{Recall}_{\text{val}}(\tau) \ge 0.80 \}
\end{equation}
This threshold is frozen before held-out test evaluation.

\section{Experimental Protocol \& Results}
Across 5 random seeds (Seeds 42, 101, 2024, 777, 9999; 951,300 persisted test predictions), models were evaluated on untouched test partitions ($N=38,052$ per seed):

\begin{table*}[t]
\centering
\caption{Multi-Seed Test Performance under Natural Prevalence ($9.42\%$) across 5 Fixed Random Seeds (Mean $\pm$ Std)}
\label{tab:main_results}
\begin{tabular}{lcccccc}
\toprule
\textbf{Model Architecture} & \textbf{PR-AUC (Mean $\pm$ Std)} & \textbf{ROC-AUC} & \textbf{Uncal. Brier} & \textbf{Recall ($\tau^*$)} & \textbf{Precision} & \textbf{Specificity} \\
\midrule
\textbf{HistGradientBoosting} & \textbf{0.3717 $\pm$ 0.0067} & \textbf{0.8483 $\pm$ 0.0026} & 0.1654 $\pm$ 0.0016 & $79.66\% \pm 0.81\%$ & \textbf{24.66\% $\pm$ 0.42\%} & \textbf{74.69\% $\pm$ 0.68\%} \\
\textbf{UniformGroupNet} & 0.3693 $\pm$ 0.0070 & 0.8478 $\pm$ 0.0026 & 0.1672 $\pm$ 0.0035 & $79.75\% \pm 0.99\%$ & 24.59\% $\pm$ 0.46\% & 74.55\% $\pm$ 0.84\% \\
\textbf{FGGN (Gated Neural)} & 0.3689 $\pm$ 0.0068 & 0.8478 $\pm$ 0.0026 & 0.1659 $\pm$ 0.0046 & $79.78\% \pm 0.61\%$ & 24.64\% $\pm$ 0.46\% & 74.61\% $\pm$ 0.75\% \\
\textbf{Logistic Regression} & 0.3673 $\pm$ 0.0063 & 0.8462 $\pm$ 0.0029 & 0.1654 $\pm$ 0.0015 & $79.60\% \pm 0.95\%$ & 24.41\% $\pm$ 0.47\% & 74.36\% $\pm$ 0.81\% \\
\textbf{EnsCVDD-Net (Paper)} & 0.3563 $\pm$ 0.0055 & 0.8442 $\pm$ 0.0030 & 0.1671 $\pm$ 0.0101 & $79.82\% \pm 1.13\%$ & 24.08\% $\pm$ 0.44\% & 73.82\% $\pm$ 0.85\% \\
\bottomrule
\end{tabular}
\end{table*}

\subsection{Key Findings}
\begin{enumerate}
    \item \textbf{Domain Grouping Gain:} Both domain-grouped models significantly outperform EnsCVDD-Net ($\Delta = +0.0126$, $p = 0.0011$).
    \item \textbf{Neutral Gating Finding:} FGGN ($0.3689$) performs identically to UniformGroupNet ($0.3693$, $p = 0.448$). Softmax gating adds parameters without ranking gain.
    \item \textbf{Novel Extensions:} `InstanceGatedNet` lifts PR-AUC to $0.3698$. An 80/20 probability ensemble of HistGBT and FGGN achieves superior ranking ($0.3728 \pm 0.0064$), proving orthogonal representation learning.
\end{enumerate}

\section{Statistical Evaluation \& Inference}
Accounting for $70\%$ training overlap across splits, we compute the Nadeau-Bengio corrected variance: $\hat{\sigma}^2_{\text{corr}} = (1/5 + 38,052/177,576)\hat{S}^2 = 0.414 \hat{S}^2$. 
\begin{itemize}
    \item \textbf{FGGN vs. UniformGroupNet:} Nadeau-Bengio $95\%$ CI $[-0.00177, +0.00095]$, $p = 0.448$; Wilcoxon $p = 0.4375$, Cohen's $d = -0.54$ (Confirms neutral finding).
    \item \textbf{FGGN vs. EnsCVDD-Net:} Nadeau-Bengio $95\%$ CI $[+0.00844, +0.01680]$, $p = 0.0011$; Cohen's $d = +5.39$ (Significant domain grouping advantage).
\end{itemize}

\section{Probability Calibration (Van Calster Hierarchy)}
Class weighting shifts logits by $\approx \ln(w_+) \approx 2.26$, degrading uncalibrated test Brier score to $0.1654$ (worse than the $0.0853$ constant-prevalence baseline $\bar{y}(1-\bar{y})$). Validation-fitted Isotonic regression reduces test Brier score to \textbf{0.0703} (an $18\%$ improvement over baseline), restoring calibration intercept to $-0.058$ and slope to $0.963$ (moderate calibration).

\section{Clinical Decision Curve Analysis (DCA)}
Evaluating Net Benefit $\text{NB} = \frac{\text{TP}}{N} - \frac{\text{FP}}{N} (\frac{p_t}{1 - p_t})$ across clinical threshold probabilities $p_t \in [1\%, 35\%]$, HistGradientBoosting and FGGN provide identical, substantial net clinical utility over ``Screen All'' and ``Screen None'' policies. At $p_t = 10\%$, FGGN yields $\text{NB} = 0.0634$ vs. $0.0581$ for EnsCVDD-Net, identifying $5.3$ additional true cases per 1,000 screened individuals at equal false alarm burden.

\section{Demographic Subgroup Dynamics}
Auditing demographic strata reveals that operating precision scales directly with underlying base prevalence: $17.74\%$ in young adults (prevalence $1.38\%$) vs. $25.71\%$ in older adults ($17.25\%$). A calibration audit shows calibration slopes remain stable ($\approx 1.0$), but intercepts shift systematically with base prevalence. Single-threshold systems must be deployed with stratum-specific awareness.

\section{Explainability \& Feature Attribution}
Learned FGGN gate allocations assign $31.8\%$ to Clinical History, $30.3\%$ to Self-Reported Health, $25.8\%$ to Lifestyle, and $12.1\%$ to Demographics. Permutation importance and Tree SHAP identify \texttt{GenHlth}, \texttt{Age}, \texttt{HighBP}, \texttt{DiffWalk}, and \texttt{Stroke} as primary drivers ($\rho = 0.847$ rank correlation between PBCC and Gini).

\section{Conclusion}
Published high-accuracy, high-precision claims on CDC BRFSS 2015 reflect artificially balanced test splits. Under a leakage-safe protocol with natural $9.42\%$ prevalence, domain grouping resolves the inductive bias mismatch of 1D convolutions ($p = 0.0011$), while learned gating provides no ranking gain over uniform pooling. An 80/20 hybrid blend of HistGBT and FGGN achieves the highest overall screening performance ($0.3728$ PR-AUC).

\bibliographystyle{IEEEtran}
\bibliography{references}

\end{document}

